\documentclass[10pt]{article}
\usepackage[margin=20mm]{geometry}
\usepackage{amsmath,amssymb,fontspec,unicode-math,xurl,hyperref}
\IfFontExistsTF{DejaVu Serif}{\setmainfont{DejaVu Serif}}{\setmainfont{Latin Modern Roman}}
\setmathfont{latinmodern-math.otf}
\setlength{\parindent}{0pt}\setlength{\parskip}{6pt}\sloppy
\title{Oracle-guided loss bounds and metric consistency in MOGIC}
\author{AI-assisted technical review}
\date{}
\begin{document}
\maketitle
Abstract: A four-label, unit-norm example disproves the label-side perturbation bound used to justify MOGIC’s oracle-guided loss: the true loss is 1.3 while the stated upper bound is 1. This invalidates a central step in the theoretical argument, independently of the observed benchmark gains. The main theorem also needs the realizability condition explicitly used in its appendix. Separate metric-definition and table-label inconsistencies limit reproducibility. Static release inspection identifies unresolved training and provenance details, without establishing that historical results used that code path. Scientific adjudication is pending.\par
\section*{Paper and inspected scope}
This report concerns Suchith Chidananda Prabhu and coauthors, MOGIC: Metadata-infused Oracle Guidance for Improved Extreme Classification, ICML 2025, PMLR 267. The 28-page proceedings PDF was read through the main method and experiments and Appendices A–H, with detailed independent checks of Appendix E’s bounds and selected numerical and metric claims. Bibliographic sources cited by the paper were not individually re-reviewed. Source pages 5, 7, 19, 20, 23 and 25 were actually viewed to check the architecture, plots, formulas and tables.\par
The public repository https://github.com/suchith720/mogic was inspected at commit `d244dff86cc4f4d18a8106b2c597a188ee278a68`. This covered its README, the full selected 001 MOGIC training entry point, the relevant configuration/import/evaluation portions of variants 002–004, a selected input-concatenation entry point and the minimal `core.py`. The actual model/loss/metric classes are imported from `xcai`; their implementation is not in this frozen repository snapshot. No author program, dependency installer, dataset or model was executed. The independently written `check\_bounds.py` uses exact rational arithmetic for the counterexample.\par
\section*{Problem 1: a false perturbation bound in the core guarantee}
Equation reference: eq:printed. Severity: Major; the theoretical derivation of oracle-guided accuracy is invalid at its stated label-side inequality. Appendix E, Lemma 2 and equations (6), (13)–(16), assert that replacing label embeddings by oracle embeddings changes the averaged triplet loss by at most \[
 \frac{2KB}{L}\sum_{l=1}^L\|z_l-z_l^*\|_2.\] Here \(g\) is \(K\)-Lipschitz, the embeddings have norm at most \(B\), and the triplet loss averages over all positive–negative pairs. The proof replaces the two class-specific incidence weights by a uniform label weight in equation (14). Those weights need not agree.\par
Counterexample satisfying the stated conditions. Take one query, four labels, one positive label (label 1), and three negative labels. Let \(e_1=(1,0)\), \(e_2=(0,1)\), \(K=B=1\), and \(g(t)=[t+3/10]_+\), using the experimental margin 0.3. Set \[x=x^*=e_1,\quad z_1^*=e_1,\quad z_1=-e_1,
 \quad z_l=z_l^*=e_2\quad(l=2,3,4).\] Every embedding has norm exactly one, and \(g\) is globally 1-Lipschitz. The oracle positive score is 1 and all negative scores are zero, so its three triplet losses are zero. With the disciple labels, the positive score is \(-1\) and each triplet loss is \(1+0.3=1.3\). Averaging three identical losses leaves \(13/10\). Only label 1 changes, by Euclidean distance 2. The paper’s claimed upper bound is therefore \[\frac{13}{10}\ \leq\ 0+\frac{2}{4}\,2=1,\] which is false. Repeating the query any number of times preserves the same averaged violation. This failure is independent of sample noise, optimization error and the omitted summation notation elsewhere in the proof. A hypothesis class can contain both displayed configurations, so restricting to oracle-realizable classes does not repair Lemma 2. ◻\par
Equation reference: eq:correct. For a query \(i\) with positive set \(S_i\) of size \(P_i\), define \(e_l=\|z_l-z_l^*\|_2\). The direct Lipschitz calculation instead gives \[
 |\ell(Z^\top x_i,y_i)-\ell(Z^{*\top}x_i,y_i)|
 \leq KB\left(\frac{\sum_{l\in S_i}e_l}{P_i}
 +\frac{\sum_{l\notin S_i}e_l}{L-P_i}\right).\] Each positive label occurs in \(L-P_i\) pairs and each negative label in \(P_i\) pairs; dividing by \(P_i(L-P_i)\) yields these coefficients. The example’s corrected upper bound is 2, which is valid. Uniform \(2/L\) coefficients agree in the balanced case \(P_i=L/2\), not for general extreme classification, where positives may be sparse.\par
Equation (7)’s query-side \(2KB\) bound follows from bounding two inner products and is not refuted by this example. Combining that valid bound with [eq:correct] (\#eq:correct) gives a valid domination of the task loss by the averaged cross loss plus \(KB\) times \[M_w=\frac1N\sum_i\|x_i-x_i^*\|_2+
 \frac1{2N}\sum_i\left(\frac{\sum_{l\in S_i}e_l}{P_i}
 +\frac{\sum_{l\notin S_i}e_l}{L-P_i}\right).\] This is a changed, class-weighted matching term, not the uniformly averaged loss in equation (5). Alternatively, a larger uniform coefficient must account for the worst class imbalance. The local inequality is repaired by this calculation; no full finite-sample theorem or empirical retraining is certified.\par
Lemma 3 explicitly averages Lemma 2’s inequalities to obtain equation (23); Theorem 5 then combines Lemmas 3 and 4, and Section 3.4 presents Theorem 1 as the central justification. Thus a false supporting lemma invalidates the supplied proof of that guarantee. This is not a claim that no suitably qualified alternative theorem exists, nor that the benchmark improvements disappear.\par
\subsection*{Realizability and the scope of the main theorem}
Appendix Lemma 3 also assumes exact oracle realizability by the disciple. Section 3.4’s standalone theorem lists boundedness, Lipschitzness and complexity control but does not state that assumption. Without it, the oracle-comparison bound is false even for binary loss. Consider a point-mass distribution with positive label \(y=1\), query embedding 1 for both models, oracle label embedding 1, and a singleton disciple label class containing only \(-1\). Use \(\ell(s,y)=(1-ys)/2\) on \([-1,1]\), so \(K=1/2\), \(B=1\) and the loss is bounded in \([0,1]\). The only disciple is an empirical minimizer of any objective over its singleton class, but its risk is 1 versus oracle risk 0. Under the usual signed empirical Rademacher definition both singleton complexities are zero. With \(N=100\) and \(\delta=0.05\), the displayed right side is \(2\sqrt{\log(20)/100}=0.346164\), and the bound fails with probability one.\par
This second example deliberately does not meet Lemma 3’s realizability assumption. It refutes only the unqualified main statement; reading that assumption into the theorem blocks this example, but leaves the independent Lemma 2 defect above. An approximation term or explicit realizability restriction is needed. Lemma 4’s brief appeal to standard Rademacher arguments does not supply a checked uniform bound for the combined loss; its constants and normalization remain unverified. The unsquared norms in equations (2)/(5), MSE terminology in Appendix C, and the experimental weights also need reconciliation before transferring an analysis to the implemented objective.\par
\section*{Problem 2: metric formulas and repeated table columns conflict}
Severity: Minor; reporting and reproducibility defects, without evidence that the historical scorer implemented the printed formulas. Appendix F.4, equations (33)/(35), puts a factor \(1/k\) in DCG but not in its ideal normalizer. Even granting a correction of the unusual printed rank-index set to an ordinary one-based predicted order, a perfect top-\(k\) list with \(k\) relevant labels gives \[\frac{k^{-1}\sum_{j=1}^k1/\log(j+1)}
 {\sum_{j=1}^k1/\log(j+1)}=\frac1k,\] not one. At \(k=5\), this caps normalized DCG at 20\%, incompatible with Table 2’s 27.44 and other listed N@5 values. The literal index definition additionally orders numeric label identifiers and can give \(i=0\), hence a zero logarithm; it should be replaced by prediction rank. The usual normalized score is \[\mathrm{nDCG}@k=
 \frac{\sum_{j=1}^k y_{\pi(j)}/\log(j+1)}
 {\sum_{j=1}^{\min(k,\|y\|_0)}1/\log(j+1)},\] where \(\pi\) orders predictions and the nonempty-relevance convention is specified. Removing \(1/k\) from the numerator, or adding the same factor to the denominator, restores the perfect-ranking normalization.\par
For the identical MOGIC(OAK)/DistilBERT row, Tables 2 and 4 report N@5=27.44 and Table 4 gives PSP@1=35.70. Table 3 labels these two numbers in the opposite columns: N@5=35.70, PSP@1=27.44. The same pattern appears in Tables 11/12 and the first dataset of Table 20. The internal conflict and original PDF layout were checked directly. It is consistent with transposed reporting columns, but raw metric outputs are needed to identify every affected row. P@1/P@5 entries are not changed by this swap. The external metric implementation was not available in the inspected snapshot, so no historical scoring bug or reversed performance ranking is asserted.\par
\section*{Empirical and source-code limits}
The six Table 2/16 P@1 gains over OAK are 0.91, 2.46, 0.59, 1.05, 0.11 and 1.69 percentage points, all positive. The headline “1–2\%” is therefore a rough characterization, not a per-dataset guarantee. Table 5’s oracle P@1 loss from ground-truth to predicted metadata is 22.54 points; its caption’s 17 differs, while the body says approximately 22. The oracle–disciple gaps in that table are 13.01 and 10.69 points, not the prose’s 8 and 6. These do not reverse the visible robustness trend. Its random replacement/removal experiments do not establish robustness to every structured corruption.\par
A training-only regularizer can leave an otherwise identical deployed architecture unchanged. Table 14, however, compares an OAK encoder at 10.30 ms against MOGIC at 14.95 ms, not a clearly matched end-to-end OAK system; that is a 45.15\% mean-latency increase for the listed objects. This cannot verify equal system latency, but neither does it disprove a zero incremental cost of guidance relative to the same complete disciple pipeline. Table 13 separately compares oracle and disciple stages; total two-stage training must include the oracle.\par
The selected released MOGIC entry points freeze teacher embeddings, supporting the stated frozen-oracle design, but pass zero weights for one cross-loss direction and label matching. They also pass `block.test.dset` as the evaluation dataset while configuring best-P@1 checkpoint selection. These are concrete release settings requiring historical run mapping and a clarified held-out selection protocol. They may be exploratory variants; imported learner semantics, original experiment commits and logs are unavailable. No claim is made that published test scores were selected by this exact code or that the full method was rerun.\par
\section*{Disposition}
Assessment | Repair verification | Author adoption | Remaining work\par
P1 Major | Weighted local bound derived | Unknown | Revised theorem and objective\par
P2 Minor | Perfect-rank normalization checked | Unknown | Metric code and table mapping\par
Source concerns: Needs Evidence | No rerun | Unknown | Historical configs and split logs\par
The core theoretical derivation needs correction. Empirical positive gains and the architectural idea of metadata-guided training survive these checks as reported observations, without independent reproduction. Scientific adjudication remains pending, with no human review recorded.\par
\section*{Evidence and review limits}
The source is the frozen PMLR proceedings PDF, identified by SHA-256\par
> `fcf38835aa7eec3d35e76bfaaccc790f37d278300134c6e45217a58e57111148`\par
Physical PDF page numbers are one based and coincide with the paper’s printed page numbers in this version. The check program was written for this review; no author program was executed. Published training runs, data collection and reported performance were not reproduced. It is not independent acceptance. No human review or author adoption is recorded. Scientific adjudication remains pending.\par
9 Suchith Chidananda Prabhu; Bhavyajeet Singh; Anshul Mittal; Siddarth Asokan; Shikhar Mohan; Deepak Saini; Yashoteja Prabhu; Lakshya Kumar; Jian Jiao; Amit Singh; Niket Tandon; Manish Gupta; Sumeet Agarwal; Manik Varma. MOGIC: Metadata-infused Oracle Guidance for Improved Extreme Classification. Proceedings of the 42nd International Conference on Machine Learning, PMLR 267 (2025). Published proceedings PDF, frozen for this review on 25 October 2026. https://proceedings.mlr.press/v267/prabhu25a.html.\par
\end{document}
